% generator_I09.tex -- Prop I.9 (Prob.): bisect a given rectilinear angle.
% Geometry and text from jemmybutton (CC-BY-SA). Build: lualatex generator_I09.tex
\documentclass[12pt]{article}
\usepackage{geometry}\geometry{a5paper, margin=1.5cm}
\usepackage[lining]{ebgaramond}
\usepackage{amssymb}
\usepackage{byrne}
\def\mpPre{u := 1cm; textLabels := false;}

\begin{document}

\defineNewPicture{%
  pair A, B, C, D, E, F;%
  A := (0, 5/3u);%
  B := (-4/3u, 0);%
  C := B xscaled -1;%
  D = whatever[B, B shifted ((C-B) rotated -60)]
    = whatever[C, C shifted ((B-C) rotated 60)];%
  E := 5/4[A, B];%
  F := 5/4[A, C];%
  byAngleDefine(B, A, D, byblue,   SOLID_SECTOR);%
  byAngleDefine(C, A, D, byyellow, SOLID_SECTOR);%
  draw byNamedAngleResized();%
  byLineDefine(B, C, byyellow, SOLID_LINE, REGULAR_WIDTH);%
  byLineDefine(A, D, byblack,  SOLID_LINE, REGULAR_WIDTH);%
  byLineDefine(D, B, byblue,   SOLID_LINE, REGULAR_WIDTH);%
  byLineDefine(C, D, byblue,   SOLID_LINE, REGULAR_WIDTH);%
  byLineDefine(A, B, byred,    SOLID_LINE, REGULAR_WIDTH);%
  byLineDefine(C, A, byred,    SOLID_LINE, REGULAR_WIDTH);%
  byLineDefine(B, E, byred,    DASHED_LINE, REGULAR_WIDTH);%
  byLineDefine(C, F, byred,    DASHED_LINE, REGULAR_WIDTH);%
  draw byNamedLine(BC,AD);%
  draw byNamedLineSeq(0)(DB,CD);%
  draw byNamedLineSeq(0)(BE,AB,CA,CF);%
  draw byLabelsOnPolygon(D, B, A, C)(ALL_LABELS, 0);%
}

% STATEMENT
\begin{center}\textbf{Book~I.\enspace Prop.\ IX. Prob.}\end{center}
\vspace{0.3cm}
\noindent
\begin{minipage}[t]{0.50\linewidth}
\itshape To bisect a given rectilinear angle
(\drawAngle{BAD,CAD}).
\end{minipage}\hfill
\begin{minipage}[t]{0.46\linewidth}
\centering\drawCurrentPicture
\end{minipage}

\vspace{0.5cm}

% PROOF
\begin{center}
Take $\drawUnitLine{AB} = \drawUnitLine{CA}$~(pr.~I.3)

draw \drawUnitLine{BC}, upon which describe
\drawLine[bottom]{CD,DB,BC}~(pr.~I.1),\\
draw \drawUnitLine{AD}.

\medskip
$\because \drawUnitLine{AB} = \drawUnitLine{CA}$~(const.)\\
and \drawUnitLine{AD} common to the two triangles\\
and $\drawUnitLine{CD} = \drawUnitLine{DB}$~(const.),

$\therefore \drawAngle{BAD} = \drawAngle{CAD}$~(pr.~I.8).
\end{center}

\begin{flushright}Q.\ E.\ D.\end{flushright}

\end{document}
